\pageTitre{CONCLUSION}

Au terme de ces \InfoDuree{} de stage au sein de \InfoEntreprise, le bilan est largement
positif. Nous avons découvert le fonctionnement d'une entreprise de lots techniques du bâtiment
et participé, au sein de l'équipe plomberie, à la réalisation des réseaux d'eau froide et
d'eau chaude sanitaire du chantier \InfoChantier. Nous y avons acquis des gestes
professionnels concrets : assemblage du PPR par polyfusion, montage et réglage des
réducteurs de pression, isolation des canalisations et essais d'étanchéité.

Sur le plan technique, l'étude de cas a montré comment les règles du NF DTU~60.11 permettent
de dimensionner une colonne montante et de justifier la pose de réducteurs de pression aux
étages bas. Elle illustre le lien direct entre la théorie vue en formation et les choix faits
sur le chantier. Ce stage a répondu à nos attentes : il nous a fait progresser en rigueur, en
autonomie et en travail d'équipe, et notre implication a été appréciée par l'encadrement.

Ce travail comporte toutefois des \textbf{limites}. Notre participation s'est concentrée sur la
distribution d'eau : nous n'avons pas travaillé sur les évacuations ni sur la mise en service
complète du surpresseur. L'étude de cas repose sur des hypothèses simplifiées et devrait être
confirmée par la note de calcul du bureau d'études. Il serait intéressant de prolonger cette
étude par l'analyse des réseaux d'évacuation et par un bilan énergétique de la production
d'eau chaude, comparant chauffe-eau électriques et pompe à chaleur.

Nous recommandons à l'entreprise de renforcer sa logistique, sa politique de sécurité et la
traçabilité de ses autocontrôles. Aux futurs stagiaires, nous recommandons de prendre des notes
chaque jour et de ne pas hésiter à poser des questions.

Enfin, ce stage a conforté notre \textbf{projet professionnel} : devenir technicien en froid et
climatisation, capable d'intervenir également sur les réseaux de fluides du bâtiment. La
proximité entre ces métiers, constatée sur le terrain, est un atout que nous comptons valoriser
dans notre insertion professionnelle.
